\section{Ecuación general}

\begin{align*}
    E+S &\Longleftrightarrow ES \longrightarrow E + P
\end{align*}

Ley de velocidad

\begin{align*}
    V &= -\frac{dS}{dt} = \frac{dP}{dt}
    \\
    \\
    V = -\frac{dS}{dt} = \frac{dES}{dt} \ &; \ V = -\frac{dES}{dt} = \frac{dP}{dt}
\end{align*}

Órdenes de reacción ($S \longrightarrow P$)

\begin{itemize}
    \item 0 - Solo de K
    \item 1 - De Sustrato y de K
\end{itemize}



Equilibrio de reacción

La reacción debe tener un equilibrio que favorezca la generación del complejo ES y luego el producto.
Para la formación del complejo ES existe:

\begin{itemize}
    \item $K_{1}$: El valor de equilibrio favorable
    \item $K_{2}$: El valor de equilibrio desfavorable
\end{itemize}

Algo a tener en cuenta es que agregar enzima cambia el orden de reacción.

\begin{align*}
    [S] < [E] &\longrightarrow \text{Rx de orden 1}
    \\
    [S] = [E] &\longrightarrow \text{Se alcanza la saturación} = V_{max}
    \\
    [S] > [E] &\longrightarrow \text{Rx de orden 0}
\end{align*}

Ecuación de Michaels-Menten:

\begin{align*}
    v &= \frac{V_{max}[S]}{K_{m} + [S]}
    \\
    Km &= \text{[S] para alcanzar }\frac{1}{2}V_{max}
\end{align*}

El valor de $K_{m}$ en la ecuación de Micahels-Menten es el conjunto de $K_{m}$ y $K_{cat}$